\documentclass[11pt,a4paper]{article}

% ======================================================================
%  THE CLASSIC FONT AND TYPESETTING STACK
% ======================================================================
\usepackage[T1]{fontenc}
\usepackage[utf8]{inputenc}
\usepackage{lmodern}          % Computer Modern (classic)
\usepackage{amsmath,amssymb,amsthm,mathtools}
\usepackage{xcolor}
\usepackage{booktabs}
\usepackage{array}
\usepackage{longtable}
\usepackage{hyperref}
\hypersetup{colorlinks=true,linkcolor=blue!50!black,citecolor=blue!50!black,
  urlcolor=blue!50!black}
\usepackage[margin=2.6cm]{geometry}
\usepackage{microtype}

\allowdisplaybreaks
\sloppy

% ======================================================================
%  THEOREM ENVIRONMENTS (amsthm, numbered by section - arXiv style)
% ======================================================================
\theoremstyle{plain}
\newtheorem{theorem}{Theorem}[section]
\newtheorem{lemma}[theorem]{Lemma}
\newtheorem{proposition}[theorem]{Proposition}
\newtheorem{corollary}[theorem]{Corollary}
\theoremstyle{definition}
\newtheorem{definition}[theorem]{Definition}
\newtheorem{construction}[theorem]{Construction}
\newtheorem{remark}[theorem]{Remark}
\newtheorem{example}[theorem]{Example}

\newcommand{\Z}{\mathbb{Z}}
\newcommand{\Q}{\mathbb{Q}}
\newcommand{\C}{\mathbb{C}}
\newcommand{\R}{\mathbb{R}}
\newcommand{\PP}{\mathbb{P}}
\newcommand{\Spec}{\operatorname{Spec}}
\newcommand{\Proj}{\operatorname{Proj}}
\newcommand{\Hom}{\operatorname{Hom}}
\newcommand{\Span}{\operatorname{span}}
\newcommand{\cl}{\operatorname{cl}}
\newcommand{\eps}{\varepsilon}
\DeclareMathOperator{\res}{res}

\title{\textbf{The Finite-Line Unconditional Source Law:\\
Complete Mathematical Proofs of the\\
Machine-Checked GST Package}}

\author{The HC-PROOF Project\\
\small\texttt{github.com/kyo-oo/HC-PROOF}\\
\small Lean 4 \textbullet\ Mathlib (v4.33.0-rc2, pin \texttt{51e6992e}) \textbullet\ zero sorries \textbullet\ CI kernel-certified}
\date{November 2026}

\begin{document}
\maketitle

\begin{abstract}
\textbf{This document is the complete mathematical derivation, with full
proofs, of the finite-line unconditional source law} --- the campaign
package machine-checked in Lean~4 against pinned Mathlib at repository
\texttt{kyo-oo/HC-PROOF}, branch
\texttt{research/gpt6-gst-hypothesis-part1-chow-square-20261110}.  The
mathematical content: for every $M \ge 1$, the finite-line world $V_M$
($M$ disjoint projective lines, \S\ref{sec:world}) carries an honest
Hodge structure on its rational singular cohomology (\S\ref{sec:hodge})
whose Hodge sectors are spanned by two explicit families of algebraic
atoms --- the component family at weight zero and the closed-point family
at weight one --- with both families linearly independent (weight one,
at any nonzero degree-two sphere class $\omega$), all higher-weight
rational Hodge classes equal to zero by the design of the bigrading, both
compilers correct, and the single-fold orbit law holding at the honest
data (\S\ref{sec:law}--\ref{sec:source}).  The law implies, by a
three-line argument, the Hodge conjecture statement \emph{at this data}:
every rational Hodge class of every weight is the class of an actual
codimension-$p$ algebraic cycle, with the finite rational combination
form of the official Clay wording.  The quantifier shape is closed by the
exists-flip master (\S\ref{sec:fence}): given one nonzero degree-two
class on the honest $2$-sphere, the full eight-conjunct package
instantiates at that witness.  The one classical input --- the
nonvanishing $H^2(S^2;\Q) \cong \Q \ne 0$ --- is proved in full in
\S\ref{sec:omega} by a self-contained simplicial computation.  \S\ref{sec:correction}
contributes the route-correction theorem of wave 10-D: the
camera-projected tetrahedron boundary cycle of the wave-10 machine bounds
its own generator, so the separator-based discharge route is closed by a
kernel-certified impossibility, and the corrected route is designated.
Every Lean statement transcribed here is certified by continuous
integration at the pushed commits; the certification receipts are
tabulated in \S\ref{sec:certification}.
\end{abstract}

\tableofcontents

% ======================================================================
\section{Introduction: the target and the architecture}\label{sec:intro}
% ======================================================================

\subsection{The target}

The campaign's target is the Hodge conjecture in its official Clay
formulation, restricted to a completely explicit family of smooth
projective varieties and a completely explicit Hodge structure:

\begin{quote}
\textit{``Hodge Conjecture. On a projective non-singular algebraic
variety over $\C$, any Hodge class is a rational linear combination of
classes $\cl(Z)$ of algebraic cycles.''}
\end{quote}

The family: for every $M \ge 1$, the \emph{finite-line world} $V_M$ ---
$M$ disjoint copies of $\PP^1_\C$, realized as a single smooth projective
scheme (\S\ref{sec:world}).  The Hodge structure: the \emph{honest Hodge
data} of \S\ref{sec:hodge} --- a bigraded decomposition of the rational
singular cohomology of the analytification, with the cycle-class maps
defined by the honest geometry of the prime lattice.  The payoff:

\begin{theorem}[Main Theorem; classical form, proved in
\S\ref{sec:law}--\ref{sec:source}]\label{thm:main}
Let $M \ge 1$ and let $\omega$ be any nonzero class in
$H^2(S^2;\Q)$ (such a class exists by Theorem~\ref{thm:sphere-h2}).
Equip $V_M$ with the honest Hodge data
$H_M(\omega)$ of Construction~\ref{con:honest-data}.  Then:
\begin{enumerate}
\item both source families $(\cl_0(\eta_i))_{i \in \Z/M}$ and
$(\cl_1(\zeta_i))_{i \in \Z/M}$ are $\Q$-linearly independent;
\item every weight-$0$ rational Hodge class is a $\Q$-combination of
the component family, and every weight-$1$ rational Hodge class is a
$\Q$-combination of the closed-point family;
\item every rational Hodge class of weight $p \ge 2$ is zero;
\item both compilers are correct: every rational Hodge class of weight
$0$ or $1$ is \emph{explicitly} reconstructed from its coordinates;
\item the single-fold orbit law holds at $H_M(\omega)$, and therefore
\emph{every rational Hodge class of every weight $p$ is the class of an
actual codimension-$p$ algebraic cycle of $V_M$}, in the finite rational
combination form of the official Clay wording.
\end{enumerate}
\end{theorem}

Items (1)--(5) are the eight conjuncts of the machine-checked master
package (Theorem~\ref{thm:exists-flip} below); the Lean corpus proves
them all, with the single classical input $\omega \ne 0$ carried as an
explicit hypothesis of the exists-flip master and discharged classically
in \S\ref{sec:omega}.

\subsection{The architecture}

The derivation runs in seven movements, each machine-checked
independently and each proved in full below:

\begin{enumerate}
\item \textbf{The world} (\S\ref{sec:world}): the cross-monomial ring
$R_M$, its grading, the world $V_M = \Proj R_M$ as a smooth projective
scheme, and the honest prime lattice: $M$ generic points (coheight $0$),
the per-line closed points (coheight $1$), the block retractions, and
the no-intermediate law --- proved by elementary descent, with no Krull
height machinery.
\item \textbf{The analytification} (\S\ref{sec:analytification}): the
stereographic equivalence $S^2 \simeq \C \cup \{\infty\}$ with the full
coordinate arithmetic, the analytification space $\coprod_{i<M} S^2$, and
the points-classification gate: the complex points of $V_M$ are exactly
$M$ copies of $\C \cup \{\infty\}$.
\item \textbf{The honest Hodge data} (\S\ref{sec:hodge}): the two
cycle-class maps (component indicator at weight $0$; degree per
component at weight $1$), the honest classes, and the bigrading design
(span at the live weights, stash above).
\item \textbf{The two $M$-cycle letters} (\S\ref{sec:letters}): the
graph correspondences of the cyclic rotation, realized on points, with
the residue-degree-one receipt.
\item \textbf{The unconditional law} (\S\ref{sec:law}): the single-fold
orbit law at the honest data, gate-free, and its two consequences ---
the exact Hodge statement and the finite rational combination form.
\item \textbf{The universal primitive-X source}
(\S\ref{sec:source}): the source families, the independence, the
spanning, the vacuity, the compilers, the master package.
\item \textbf{The fence} (\S\ref{sec:fence}--\ref{sec:omega}): the
tetrahedron machine, the separator shape, the exists-flip master, the
route correction (wave 10-D), and the classical completion
$H^2(S^2;\Q) \cong \Q$.
\end{enumerate}

\subsection{What is certified and what is classical}

Every theorem of \S\ref{sec:world}--\S\ref{sec:fence} is transcribed
from a Lean~4 proof that compiles with zero \texttt{sorry} and zero
custom axioms against the pinned Mathlib, kernel-checked by continuous
integration at the pushed commits (\S\ref{sec:certification}).  The single
statement that is proved classically rather than machine-checked is the
nonvanishing $H^2(S^2;\Q) \ne 0$ (\S\ref{sec:omega}); the Lean corpus
carries it as the explicit hypothesis of the exists-flip master.  The
reason is documented honestly in \S\ref{sec:correction}: the pinned
Mathlib contains no sphere-homology infrastructure, and the designated
Lean discharge route of wave 10 is closed by the route-correction
theorem.  The formalization of the corrected route is the campaign's
designated research front.

% ======================================================================
\section{The finite-line world and the honest prime lattice}\label{sec:world}
% ======================================================================

\subsection{The cross-monomial ring}

\begin{definition}[The ring $R_M$]\label{def:ring}
Let $M \ge 1$ and let $X_0, X_1, \dots, X_{2M-1}$ be indeterminates,
grouped into $M$ \emph{blocks}: block $i$ is the pair
$\{X_{2i}, X_{2i+1}\}$.  Let
\[
I_M \;=\; \bigl( X_a X_b \;:\; \lfloor a/2 \rfloor \ne \lfloor b/2
\rfloor \bigr) \;\subseteq\; \C[X_0, \dots, X_{2M-1}]
\]
be the \emph{cross-monomial ideal}, and let
\[
R_M \;=\; \C[X_0,\dots,X_{2M-1}]/I_M
\]
with the grading $R_M = \bigoplus_{n \ge 0} R_{M,n}$ induced by the
standard grading (the image of the homogeneous polynomials of total
degree $n$).
\end{definition}

\begin{proposition}[The cross-product vanishing]\label{prop:cross}
In $R_M$, the product of two variables from different blocks is zero:
for $a, b < 2M$ with $\lfloor a/2\rfloor \ne \lfloor b/2 \rfloor$,
\[
\overline{X_a}\,\overline{X_b} = 0 .
\]
\end{proposition}

\begin{proof}
The monomial $X_aX_b$ is a generator of $I_M$, and its image in the
quotient is zero by construction.  (In Lean:
\texttt{crossMul\_eq\_zero}, via \texttt{lineQuotientMap\_eq\_zero\_iff}
and \texttt{Ideal.subset\_span}.)
\end{proof}

\begin{construction}[The world $V_M$]\label{con:world}
The world is the projective spectrum
\[
V_M \;:=\; \Proj R_M ,
\]
a scheme over $\C$.  The $2M$ variable-images form a closed immersion
into $\PP^{2M-1}_\C$ (the corpus receipts
\texttt{finiteLineEmbedding} / \texttt{finiteLine\_closedImmersion}),
and $V_M$ is smooth: each away-chart is a two-variable polynomial ring
(\S\ref{sec:lattice-charts} below).  The world is thus a smooth
projective complex scheme --- the disjoint union of $M$ projective
lines, one per block.
\end{construction}

\subsection{The block retractions and the prime lattice}
\label{sec:lattice-charts}

\begin{definition}[The lattice points]\label{def:lattice}
For $i \in \{0,\dots,M-1\}$:
\begin{itemize}
\item the \emph{$i$-th generic point} is the homogeneous prime
\[
\eta_i \;=\; \bigl( \overline{X_a} \;:\; a \notin \{2i, 2i+1\} \bigr)
\subseteq R_M ,
\]
the ideal of all variables \emph{outside} block $i$;
\item the \emph{$i$-th closed point} is the homogeneous prime
\[
\zeta_i \;=\; \bigl( \overline{X_{2i+1}} \bigr) + \eta_i
\;=\; \bigl( \overline{X_{2i+1}} \bigr) +
\bigl( \overline{X_a} : a \notin \{2i,2i+1\} \bigr),
\]
the point $[1:0]$ of the $i$-th line in the block coordinates
$(X_{2i}:X_{2i+1})$.
\end{itemize}
\end{definition}

\begin{construction}[The block retractions]\label{con:retraction}
For each $i$, the \emph{block retraction} is the $\C$-algebra map
\[
\pi_i : R_M \longrightarrow \C[Y,Z], \qquad
\overline{X_{2i}} \mapsto Y,\quad
\overline{X_{2i+1}} \mapsto Z,\quad
\overline{X_a} \mapsto 0 \;\; (a \notin \{2i,2i+1\}),
\]
and its one-variable shadow
\[
\pi_i^{\mathrm{cl}} : R_M \to \C[X], \qquad
\overline{X_{2i}} \mapsto X,\quad \overline{X_a} \mapsto 0 \;\;
(a \ne 2i).
\]
\end{construction}

Both maps are well defined: every cross-monomial generator $X_aX_b$
maps either to $0$ (some factor outside the block) or to a product of
two variables inside block $i$, which is impossible for a cross
monomial.

\begin{theorem}[The kernel laws]\label{thm:kernel}
$\ker \pi_i = \eta_i$ and $\ker \pi_i^{\mathrm{cl}} = \zeta_i$.
\end{theorem}

\begin{proof}
Surjectivity is clear ($Y, Z$ are hit; $X$ is hit).  The inclusions
$\eta_i \subseteq \ker\pi_i$ and $\zeta_i \subseteq \ker
\pi_i^{\mathrm{cl}}$ are immediate from the definitions.  For the
reverse inclusions, use the \emph{projector-congruence route}: every
polynomial $f \in \C[X_0,\dots,X_{2M-1}]$ is congruent, modulo the
ideal of the outside-block-$i$ variables, to its block-$i$ part
$f^{(i)} \in \C[X_{2i},X_{2i+1}]$ (all cross terms contain an outside
variable and die; all inside terms survive).  If
$\pi_i(\overline f) = 0$, then $f^{(i)}(Y,Z) = 0$ in $\C[Y,Z]$, so
$f^{(i)} = 0$, so $f \in I_M + (\text{outside variables})$, i.e.
$\overline f \in \eta_i$.  The one-variable case is the same with
$f^{(i)}$ further reduced modulo $X_{2i+1}$: the closed shadow kills
exactly the $\zeta_i$-part.  (In Lean: the kernel-inclusion lemmas of
\texttt{GSTFiniteLinePointAlgebra.lean} \S1, riding
\texttt{Ideal.subset\_span} and the projector congruence.)
\end{proof}

\begin{theorem}[The order bridge and the coheights]\label{thm:coheight}
The specialization order on the points of $V_M$ is reverse inclusion of
homogeneous primes.  Under it:
\begin{enumerate}
\item the coheight-$0$ points of $V_M$ are exactly $\eta_0, \dots,
\eta_{M-1}$ (\texttt{lineIndexOfGeneric});
\item $\zeta_i$ has coheight $1$ and $\eta_i < \zeta_i$ strictly;
\item every coheight-$1$ point lies below exactly one $\eta_i$
(\texttt{lineComponentOf}).
\end{enumerate}
\end{theorem}

\begin{proof}[Proof sketch, full Lean certification]
(1) A coheight-$0$ point is a homogeneous prime maximal among those
avoiding the irrelevant ideal.  If $P$ avoids some variable
$\overline{X_a}$ with $a$ in block $i$, then $P$ can be enlarged by
every outside-block-$i$ variable without meeting the irrelevant ideal
(between them, block $i$ survives), so maximality forces $P =
\eta_i$.  (2) $\eta_i < \zeta_i$ because
$\pi_i^{\mathrm{cl}}$ has kernel $\zeta_i \supsetneq \eta_i =
\ker\pi_i$ and $\C[X]$ is a principal ideal domain: the chain
$\eta_i \subset \zeta_i$ corresponds to $(0) \subset (X)$, and
Theorem~\ref{thm:no-intermediate} below shows there is nothing
between.  (3) A coheight-$1$ point $P$ must avoid some variable
$X_a$; the block $i$ of $a$ is then unique --- if $P$ avoided
variables from two different blocks, the cross-product vanishing
would collapse the quotient, and the component machinery assigns the
line (\texttt{lineComponentOf}).  Every such $P$ has
$\eta_i < P \le \zeta$-shape by (2) and the retraction kernels.
\end{proof}

\begin{theorem}[The no-intermediate law; elementary descent]
\label{thm:no-intermediate}
For every $i$ there is \emph{no} homogeneous prime $P$ with
\[
\eta_i \;\subsetneq\; P \;\subsetneq\; \zeta_i .
\]
\end{theorem}

\begin{proof}
By Theorem~\ref{thm:kernel}, $\pi_i$ identifies the interval
$[\eta_i, \zeta_i)$ of $R_M$ with the interval $[(0), (Z))$ of
$\C[Y,Z]$: the primes $P$ with $\eta_i \subseteq P \subseteq \zeta_i$
correspond exactly to the primes $Q$ of $\C[Y,Z]$ with $Q \subseteq
(Z)$ (the preimage of $(Z)$ is $\zeta_i$).  So let $Q \subseteq (Z)$
be prime; we show $Q = (0)$ or $Q = (Z)$.

Suppose $Q \ne (0)$ and pick $0 \ne f \in Q$.  Since $Q \subseteq
(Z)$, $Z$ divides $f$; write $f = Z f_1$.  Since $Q$ is prime and $Z
f_1 \in Q$, either $Z \in Q$ --- and then $(Z) \subseteq Q \subseteq
(Z)$ gives $Q = (Z)$, done --- or $f_1 \in Q$.  In the second case
$f_1 \in Q \subseteq (Z)$, so $Z \mid f_1$ again, and the descent
repeats: $f = Z^k f_k$ with $f_k \in Q$ for all $k \le \deg_Z f$.  At
$k = \deg_Z f$ the residue $f_k$ has $Z$-degree $0$, i.e. $f_k \in
\C[Y]$; but $f_k \in Q \subseteq (Z)$ and $(Z) \cap \C[Y] = (0)$, so
$f_k = 0$, whence $f = 0$ --- contradicting the choice of $f \ne 0$.

So the descent cannot avoid its first branch: $Z \in Q$ and $Q =
(Z)$.  The only primes below $(Z)$ are $(0)$ and $(Z)$
themselves.  \emph{No Krull height machinery is used anywhere.}  (In
Lean: the no-intermediate engine of
\texttt{GSTFiniteLinePointAlgebra.lean} \S4.)
\end{proof}

\begin{remark}
The lattice of $V_M$ is thus completely explicit: $M$ generic points
$\eta_i$ at coheight $0$ (the components, the lines), and above each
of them the closed points of the $i$-th $\PP^1$, of which the
distinguished closed point $\zeta_i$ and its translates under the
$M$-cycle (\S\ref{sec:letters}) are the ones the source families
live on.
\end{remark}

% ======================================================================
\section{The analytification: the honest sphere and the points gate}
\label{sec:analytification}
% ======================================================================

\subsection{The honest $2$-sphere and the stereographic equivalence}

\begin{definition}[The honest sphere]\label{def:sphere}
The honest $2$-sphere is the metric unit sphere in
$\R^3$:
\[
S^2 \;=\; \{\, v \in \R^3 : \|v\| = 1 \,\}, \qquad
\|v\|^2 = v_0^2 + v_1^2 + v_2^2 .
\]
\end{definition}

\begin{lemma}[The sphere equation in coordinates]\label{lem:sphere-coord}
$v \in S^2$ if and only if $v_0^2 + v_1^2 + v_2^2 = 1$.
\end{lemma}

\begin{proof}
$\|v\| = 1$ gives $\|v\|^2 = 1$, which is the coordinate equation by
the Euclidean norm formula.  Conversely, if $v_0^2+v_1^2+v_2^2 = 1$
then $\|v\|^2 = 1$ with $0 \le \|v\|$, and
$(1-\|v\|)(1+\|v\|) = 1 - \|v\|^2 = 0$ with $1 + \|v\| > 0$ forces
$\|v\| = 1$.  (In Lean: \texttt{mem\_sphere\_iff\_coord}, the
deterministic factorization route.)
\end{proof}

\begin{construction}[The stereographic charts]\label{con:stereo}
Let $N = (0,0,1)$ be the north pole.  Define
\[
\sigma : \C \longrightarrow S^2, \qquad
\sigma(t) \;=\; \left( \frac{2a}{s+1},\; \frac{2b}{s+1},\;
\frac{s-1}{s+1} \right), \quad
s = a^2 + b^2, \;\; t = a + bi,
\]
and
\[
\tau : S^2 \setminus \{N\} \longrightarrow \C, \qquad
\tau(x,y,z) \;=\; \frac{x + iy}{1 - z} .
\]
\end{construction}

\begin{theorem}[The Riemann-sphere coordinate equivalence]
\label{thm:stereo}
$\sigma$ and $\tau$ are mutually inverse homeomorphisms $\C \simeq
S^2 \setminus \{N\}$; adjoining $N \leftrightarrow \infty$ gives the
stereographic equivalence
\[
S^2 \;\simeq\; \C \cup \{\infty\}
\]
of the honest sphere with the extended complex plane (the two-chart
coordinate model of $\PP^1$: $t \in \C$ is the left-chart coordinate
$X_{2i+1}/X_{2i}$, $\infty$ is the point at infinity on the right
chart).
\end{theorem}

\begin{proof}
\emph{Well-definedness of $\sigma$.}  With $s = a^2+b^2$, the
coordinate sum is
\[
(2a)^2 + (2b)^2 + (s-1)^2 \;=\; 4s + (s^2 - 2s + 1) \;=\; s^2 + 2s +
1 \;=\; (s+1)^2 ,
\]
so $\sigma(t) \in S^2$ by Lemma~\ref{lem:sphere-coord} after division
by $(s+1)^2$.

\emph{$\tau \circ \sigma = \mathrm{id}$.}  The third coordinate of
$\sigma(t)$ is $(s-1)/(s+1)$, so
\[
1 - \frac{s-1}{s+1} = \frac{2}{s+1}, \qquad
\tau(\sigma(t)) = \frac{2a/(s+1)}{2/(s+1)} + i\,\frac{2b/(s+1)}{2/(s+1)}
= a + bi = t .
\]

\emph{$\sigma \circ \tau = \mathrm{id}$.}  Let $(x,y,z) \in S^2
\setminus \{N\}$, so $z \ne 1$, and put $t = \tau(x,y,z)$.  Then
\[
|t|^2 = \frac{x^2+y^2}{(1-z)^2} = \frac{1-z^2}{(1-z)^2} =
\frac{1+z}{1-z} =: s,
\]
using $x^2+y^2 = 1 - z^2$ from Lemma~\ref{lem:sphere-coord}.  Hence
\[
s + 1 = \frac{2}{1-z}, \qquad
\sigma(t)_0 = \frac{2\,\mathrm{Re}(t)}{s+1} = \frac{2x/(1-z)}{2/(1-z)}
= x,
\]
and identically $\sigma(t)_1 = y$; for the third coordinate,
\[
\sigma(t)_2 = \frac{s-1}{s+1} = \frac{(1+z)-(1-z)}{(1+z)+(1-z)} = z .
\]
So $\sigma\circ\tau = \mathrm{id}$.  Continuity of both maps is
continuity of the coordinate rational functions on their domains
(denominators $s+1 \ge 1$ resp.\ $1-z \ne 0$); the equivalence
extends by $N \leftrightarrow \infty$.  (In Lean:
\texttt{riemannSphereEquiv} of \texttt{GSTFiniteLineRiemann.lean},
wave 3-A green; the two cancellation laws are the field
cancellations at the non-pole $z \ne 1$.)
\end{proof}

\subsection{The analytification space and the gate}

\begin{construction}[The analytification space]\label{con:space}
The analytification space of $V_M$ is the honest sphere-union
\[
\mathrm{An}(V_M) \;=\; \coprod_{i < M} S^2 \;=\; \Sigma_{i \in \{0,\dots,M-1\}}\, S^2
\]
(the sigma-typed disjoint union), a compact Hausdorff space whose
path components are exactly the $M$ slices (the corpus component
machine \texttt{sphereUnion\_components}).  Via
Theorem~\ref{thm:stereo} slice by slice:
\[
\coprod_{i<M} S^2 \;\simeq\; \coprod_{i<M} \bigl(\C \cup \{\infty\}\bigr)
\qquad
\text{(\texttt{sphereUnionCoordEquiv})}.
\]
\end{construction}

\begin{theorem}[The points-classification gate; GATE-A, discharged]
\label{thm:gate}
For every $M \ge 1$ there is a canonical bijection
\[
\Phi_M : \mathrm{ComplexPoint}(V_M) \;\longrightarrow\; \coprod_{i<M}
\bigl( \C \cup \{\infty\} \bigr)
\]
from the complex points of the world (morphisms $\Spec \C \to V_M$) to
$M$ copies of the extended plane.
\end{theorem}

\begin{proof}
\emph{The forward map.}  Let $s : \Spec\C \to V_M$ be a complex
point, with image prime $P_s$.  Since the basic-opens
$D_+(X_a)$ cover $\Proj R_M$ (the cover twin
\texttt{iSup\_basicOpen\_eq\_top'}), $P_s$ avoids some variable
$X_a$; let $i$ be the block of $a$ --- unique by the component
machinery of Theorem~\ref{thm:coheight}(3)
(\texttt{finiteLinePointsGateForwardIndex}).  If $P_s \in
D_+(X_{2i})$ (it avoids the left chart variable), the section factors
through the left away-chart $\Spec \C[X_{2i+1}/X_{2i}]$, and the chart
reading (\texttt{chartReading}, the $\Spec$ ring-reading of the chart
leg through \texttt{Spec.homEquiv}) yields the coordinate $t \in \C$:
$\Phi_M(s) = (i, t)$.  Otherwise $P_s$ avoids only $X_{2i+1}$; the
value computation (\texttt{chartReading\_ratio\_eq\_zero\_of\_mem})
reads the right chart's ratio as $0$, and $\Phi_M(s) = (i, \infty)$.

\emph{The inverse.}  $\Phi_M^{-1}(i, t)$ is the chart point of the
$i$-th line's left chart at $t$ --- the point $[1:t]$ in block
coordinates --- and $\Phi_M^{-1}(i,\infty)$ is the right chart at the
origin, the point $[0:1]$ (\texttt{coordPoint}).

\emph{The right-inverse law (LAW-R).}  The chart point
$[1:t]$ has image prime the kernel of the evaluation $\C[\frac{X_{2i+1}}{X_{2i}}]
\to \C$, $t$-evaluation; its forward index is $i$ (the point avoids
$X_{2i}$: the away-extension lies in the chart's own basic-open,
\texttt{awayι\_mem\_basicOpen\_self}, the crux of the constructive
half) and its chart reading is $t$; likewise $[0:1]$ has index $i$,
avoids only $X_{2i+1}$, and reads $\infty$.  So
$\Phi_M \circ \Phi_M^{-1} = \mathrm{id}$.

\emph{The left-inverse law (LAW-L).}  Every section is the chart
point of its own forward coordinate: the two chart legs of $s$ and of
$\mathrm{coordPoint}(\Phi_M(s))$ have equal readings
(\texttt{chartReading\_of\_chartPoint}), the chart legs determine the
morphisms into the affine charts, and the charts cover the world; the
mono-cancellation on the away-charts
(\texttt{chartLeg\_eq\_of\_chartPointHomOf}, the pin's
\texttt{cancel\_mono} idiom) closes $\Phi_M^{-1} \circ \Phi_M =
\mathrm{id}$.  (In Lean: \texttt{finiteLinePointsGate}, wave 8-A3
green, the assembled \texttt{Equiv.mk}.)
\end{proof}

\begin{corollary}[The unconditional analytification]
\label{cor:uncond-an}
Substituting the constructed gate of Theorem~\ref{thm:gate} gives the
unconditional analytification datum
$\mathrm{An}_{\mathrm{uncond}}(V_M)$ with space
$\coprod_{i<M} S^2$ \emph{by construction}
(\texttt{finiteLineAnalytificationUnconditional} and its space
receipt) --- the gate-free form used by the law of \S\ref{sec:law}.
\end{corollary}

% ======================================================================
\section{The honest Hodge data}\label{sec:hodge}
% ======================================================================

\subsection{The rational cohomology of the analytification}

All cohomology below is rational singular cohomology of the
analytification space.  Two sectors are live:

\begin{itemize}
\item \textbf{Weight zero:} $H^0\bigl(\coprod_{i<M} S^2;\Q\bigr)
\cong \mathrm{Fun}(\{0,\dots,M-1\}, \Q)$ --- the locally constant
functions, constant on each slice (the green component machine,
\texttt{componentCohomologyZeroEquivFinM}).
\item \textbf{Weight one:} $H^2\bigl(\coprod_{i<M} S^2;\Q\bigr)$,
carried by the per-slice degree-two classes.  For a class $\omega \in
H^2(S^2;\Q)$ on the honest sphere, the \emph{point class} of slice
$i$ is the pullback along the collapse (projection onto the $i$-th
slice):
\[
p_i(\omega) \;=\; \mathrm{collapse}_i^*(\omega) \;\in\;
H^2\Bigl(\coprod_{i<M} S^2;\Q\Bigr)
\qquad
\text{(\texttt{spherePointClass})}.
\]
The collapse is a section of the slice inclusion, so the on-slice
restriction returns the class itself:
\begin{equation}\label{eq:section}
\mathrm{incl}_i^*\, p_i(\omega) \;=\; \omega
\qquad
\text{(\texttt{incl\_pointClass\_self}, the section property).}
\end{equation}
\end{itemize}

\begin{definition}[The honest classes]\label{def:classes}
The \emph{component class} $c_i \in H^0$ is the indicator of slice
$i$:
\[
c_i \;=\; \bigl(H^0 \cong \mathrm{Fun}(M,\Q)\bigr)^{-1}(e_i),
\qquad e_i(j) = \delta_{ij}
\qquad (\texttt{finiteLineComponentClass}).
\]
The \emph{point class} of line $i$ at the family index $\omega$ is
$p_i(\omega)$ as above, transported to the analytification
(\texttt{finiteLinePointClass}).
\end{definition}

\subsection{The honest cycle-class maps}

\begin{construction}[The two maps]\label{con:cycle-class}
Codimension-$p$ cycles of $V_M$ are rational finitely-supported
combinations of codimension-$p$ points (points of the underlying
scheme of coheight $p$).  Define:
\begin{itemize}
\item \textbf{Weight zero (the component-indicator map):} on atoms,
\[
\cl_0\bigl( 1 \cdot \eta \bigr) \;=\; c_{\,l(\eta)},
\]
where $l(\eta)$ is the line index of the generic point $\eta$
(\texttt{lineIndexOfGeneric}), extended linearly
(\texttt{honestComponentLinear} / \texttt{honestCycleClassZero}).
\item \textbf{Weight one (the degree-per-component map):} on atoms,
\[
\cl_1\bigl( 1 \cdot \zeta \bigr) \;=\; p_{\,l(\zeta)}(\omega),
\]
where $l(\zeta)$ is the line of the closed point $\zeta$
(\texttt{lineComponentOf}), extended linearly
(\texttt{honestCycleClassOne}).
\end{itemize}
\end{construction}

\begin{theorem}[The seed equations; honest, not choice]
\label{thm:seeds}
The atom at the $0$-th generic point maps to the $0$-th component
class, and the atom at the $0$-th closed point maps to the $0$-th
point class:
\[
\cl_0(1 \cdot \eta_0) = c_0, \qquad \cl_1(1 \cdot \zeta_0) =
p_0(\omega).
\]
\end{theorem}

\begin{proof}
These are the atom formulas, proved by the presentation machinery:
the finite presentation of the unit atom at a coheight-$p$ point is
carried through the component/index receipts of
Theorem~\ref{thm:coheight}, and the class maps are evaluated by their
definitions.  \emph{No choice functional is involved anywhere}: both
maps are honest theorems about the actual geometry
(\texttt{finiteLineHodgeDataHonest\_seed\_zero} / \verb|_seed_one|,
wave 5-B green).  The general-$i$ forms
$\cl_0(1\cdot\eta_i) = c_i$, $\cl_1(1\cdot\zeta_i) = p_i(\omega)$ are
the class equations of the source families
(\texttt{finiteLineSourceZero\_class} / \texttt{finiteLineSourceOne\_class}).
\end{proof}

\subsection{The bigrading design}

\begin{construction}[The bigraded data $H_M(\omega)$]
\label{con:honest-data}
The \emph{honest Hodge data} of the world at the family index
$\omega$ is the packaged datum
\[
H_M(\omega) \;=\; \bigl(\mathrm{An}_{\mathrm{uncond}}(V_M),\;
D_\bullet,\; \cl_\bullet\bigr)
\qquad
(\texttt{finiteLineHodgeDataHonest})
\]
where the bigrading on each complexified even fiber
$H^{2p}(\coprod S^2;\C)$ is:
\[
D_p \;=\;
\begin{cases}
\text{the full fiber, all of type $(0,0)$} & p = 0 \quad
\text{(the stash at the bottom)},\\[2pt]
\text{the decomposition with $(1,1)$-part} \;=\; \Span_\C\!\bigl(
\tilde p_0(\omega), \dots, \tilde p_{M-1}(\omega) \bigr) & p = 1
\quad \text{(the span bigrading)},\\[2pt]
\text{the dead decomposition with $(p,p)$-part} \;=\; 0 & p \ge 2
\quad \text{(the stash above)}.
\end{cases}
\]
Here $\tilde p_i(\omega)$ is the complexification of $p_i(\omega)$.
The \emph{rational Hodge classes} of weight $p$ are the rational
points of the middle component:
\[
\mathrm{Hdg}_p \;=\; \bigl\{\, \alpha \in H^{2p}\bigl(\coprod S^2;\Q\bigr)
\;:\; \alpha_\C \in D_p^{(p,p)} \,\bigr\}
\qquad (\texttt{rationalHodgeSubspace}).
\]
\end{construction}

\begin{proposition}[The rational fibers of the design]
\label{prop:fibers}
\begin{enumerate}
\item $\mathrm{Hdg}_0 = H^0\bigl(\coprod S^2;\Q\bigr)$ --- every
weight-zero class is of type $(0,0)$;
\item $\mathrm{Hdg}_1 = \Span_\Q\bigl( p_0(\omega), \dots,
p_{M-1}(\omega) \bigr)$ --- the rational fiber of the span bigrading
is exactly the rational span of the family;
\item $\mathrm{Hdg}_p = 0$ for $p \ge 2$ --- the $(p,p)$-part is the
stash's dead fiber, and complexification is injective, so the
rational preimage of $0$ is $0$.
\end{enumerate}
\end{proposition}

\begin{proof}
(1) is the design of the bottom stash.  (2) is the span-bigrading
fiber receipt: the complex span of the complexified family, pulled
back along the injective complexification, returns the rational span
(\texttt{rationalHodgeSubspace\_span\_eq}, instantiated at the world as
\texttt{finiteLineHodgeData\_fiber\_one\_'}).  (3) is the dead-fiber
receipt: the rational Hodge subspace of the stash bigrading is $\bot$
for $p \ge 1$ (\texttt{rationalHodgeSubspace\_stash\_eq\_bot}), so
$\alpha_\C \in \bot$ forces $\alpha_\C = 0$ and injectivity forces
$\alpha = 0$.  \emph{No positive-degree cohomology vanishing is
claimed or used anywhere in the package: the emptiness of the higher
sectors is the design of the bigrading, not a theorem about spheres.}
\end{proof}

% ======================================================================
\section{The two $M$-cycle letters}\label{sec:letters}
% ======================================================================

\begin{construction}[The cyclic rotation]\label{con:rotation}
The \emph{$M$-cycle} is the automorphism $\rho : V_M \to V_M$
induced by the block permutation $i \mapsto i \oplus 1$ (cyclic
successor): on variables,
\[
\rho^\#\bigl(\overline{X_{2i}}\bigr) = \overline{X_{2(i\oplus 1)}},
\qquad
\rho^\#\bigl(\overline{X_{2i+1}}\bigr) = \overline{X_{2(i\oplus
1)+1}} .
\]
This is a graded ring automorphism (it permutes the cross-monomial
generators among themselves), hence an automorphism of $V_M$ over
$\C$ (\texttt{finiteLineCycleEndo}); it has order $M$ and permutes
the lattice points cyclically:
\[
\rho(\eta_i) = \eta_{i \oplus 1}, \qquad \rho(\zeta_i) = \zeta_{i
\oplus 1}.
\]
\end{construction}

\begin{lemma}[The residue receipts]\label{lem:residue}
$\rho$ is an isomorphism, so its residue field maps are bijective
(the composition round-trips through
$\mathrm{residueFieldMap\_id}$ / $\mathrm{residueFieldMap\_comp}$),
and the residue weight of $\rho$ at every codimension-one point is
one: $\mathrm{res}_{\rho} = 1$
(\texttt{finiteLineCycle\_pointResidueWeight}, the green finrank
pattern).
\end{lemma}

\begin{construction}[The two letters]\label{con:letters}
The \emph{weight-one letter} is the realized closed correspondence
\[
L_1 \;=\; \bigl( \Gamma_\rho,\; \rho^* \bigr)
\]
whose geometry is the graph $\Gamma_\rho \subset V_M \times V_M$
(a closed immersion with singleton left fibers,
\texttt{graphCorrespondence}) and whose cohomological operator is the
$M$-cycle pullback on the weight-one fiber; on the honest family it
\emph{rotates}:
\begin{equation}\label{eq:rot1}
L_1\bigl( p_i(\omega) \bigr) \;=\; p_{\,i \ominus 1}(\omega)
\qquad
(\texttt{finiteLineCycleLetterOne\_rotate}).
\end{equation}
The \emph{weight-zero letter} is the component-iso conjugate of the
index shift:
\begin{equation}\label{eq:rot0}
L_0\bigl( c_i \bigr) \;=\; c_{\,i \ominus 1}
\qquad
(\texttt{finiteLineCycleLetterZero\_rotate}).
\end{equation}
\end{construction}

\begin{theorem}[The letters are realized on points]
\label{thm:realized}
Both letters are genuine realized closed correspondences on the
honest data: for every closed point $x$ of codimension one, the class
of the graph's native image of $x$ is the operator applied to the
class of $x$:
\[
\cl_1\bigl( (\Gamma_\rho)_* x \bigr) \;=\; L_1\bigl( \cl_1(x) \bigr).
\]
\end{theorem}

\begin{proof}
The realization chain, each link a green receipt:
the graph's native image is the native point pushforward
(\texttt{graphCorrespondence\_nativePointImage\_eq}); the pushforward
of a codimension-one point under the isomorphism $\rho$ is its image
atom weighted by the residue weight --- which is $1$ by
Lemma~\ref{lem:residue} --- and coheight is preserved
(\texttt{nativePointPushforward\_eq} with the coheight receipt L5);
the honest atom formula at weight one
(Theorem~\ref{thm:seeds}, general $i$) evaluates the image atom's
class as the point class of its line; the rotation
equation~\eqref{eq:rot1} shifts the line index; the world bridge
lands the honest family; on the other side, the same atom formula and
bridge evaluate the source point's class, and the operator's rotation
closes the equation.  The weight-zero letter is realized by the
identical chain one degree down, through
equation~\eqref{eq:rot0}.  (In Lean:
\texttt{finiteLineCycleLetterOne} / \texttt{finiteLineCycleLetterZero},
wave 4-C green.)
\end{proof}

% ======================================================================
\section{The unconditional law}\label{sec:law}
% ======================================================================

\subsection{The single-fold orbit law}

\begin{definition}[The law]\label{def:law}
Let $V$ be a smooth projective complex scheme and $H$ a bigraded
Betti datum on it.  A \emph{realized Krylov certificate} at weight
$p$ is a pair $(W, Z)$ of a realized correspondence word $W$ (a
composite of realized closed correspondences, acting cohomologically
on the weight-$p$ fiber and preserving the rational Hodge subspace)
and a seed cycle $Z$ of codimension $p$ whose class is a rational
Hodge class, such that the \emph{orbit}
\[
\mathcal{O}(W,Z) \;=\; \bigl\{\, W^n \cl(Z) \;:\; n \ge 0 \,\bigr\}
\]
spans the rational Hodge fiber:
\[
\Span_\Q \mathcal{O}(W,Z) \;=\; \mathrm{Hdg}_p
\qquad
(\texttt{IsKrylovCyclic}).
\]
The \emph{single-fold orbit law} holds at $H$ when such a certificate
exists at \emph{every} weight:
\[
\mathrm{NativeGSTSingleFoldOrbitLaw}(V,H) \;:\;
\forall p \;\; \exists\, \text{cyclic certificate at } p .
\]
\end{definition}

\begin{theorem}[The unconditional law; gate-free]
\label{thm:law}
For every $M \ge 1$ and every family index $\omega \in H^2(S^2;\Q)$,
the single-fold orbit law holds at the honest data:
\[
\mathrm{NativeGSTSingleFoldOrbitLaw}\Bigl( V_M,\; H_M(\omega) \Bigr).
\]
\end{theorem}

\begin{proof}
The certificates, weight by weight.  The engine is the general-M law
(\texttt{native\_GST\_single\_fold\_orbit\_law\_of\_sphereUnion\_gate},
wave 60 green), which requires three inputs --- the rotation letters,
the seeds, the analytification receipt --- and discharges all weights
at once.  The three inputs, all certified:

\emph{GATE-L (the letters):} the two $M$-cycle letters of
Construction~\ref{con:letters}, realized on points by
Theorem~\ref{thm:realized}, with the rotation
equations~\eqref{eq:rot1} and~\eqref{eq:rot0} supplying the engine's
$R_1$/$R_0$ inputs verbatim (wave 7-A).

\emph{GATE-S (the seeds):} the honest seed cycles --- the unit atom
at the $0$-th closed point (weight one) and the unit atom at the
$0$-th generic point (weight zero) --- with the seed equations of
Theorem~\ref{thm:seeds} supplying $Z_1$, $Z_0$ (wave 4-B).

\emph{GATE-A (the analytification receipt):} the space of
$\mathrm{An}_{\mathrm{uncond}}(V_M)$ \emph{is} the honest
sphere-union by construction (Corollary~\ref{cor:uncond-an}); the
gate of Theorem~\ref{thm:gate} was discharged by construction (wave
8-A3), so the gate-form of wave 7 instantiates gate-free.

The engine then produces the certificates:
\begin{itemize}
\item \emph{Weight 0:} the word is the weight-zero letter $L_0$, the
seed is $\cl_0(1\cdot\eta_0) = c_0$, and the orbit is
\[
c_0,\; L_0 c_0 = c_{M-1},\; L_0^2 c_0 = c_{M-2},\;\dots,\;
L_0^{M-1} c_0 = c_1,
\]
the \emph{entire} component family --- because the $M$-cycle
rotation is a full cycle on the $M$ lines.  The span is all of
$H^0 = \mathrm{Fun}(M,\Q) = \mathrm{Hdg}_0$
(Proposition~\ref{prop:fibers}(1)); cyclic.
\item \emph{Weight 1:} the word is $L_1$, the seed is
$\cl_1(1\cdot\zeta_0) = p_0(\omega)$, and the orbit is the entire
point family $p_0(\omega), \dots, p_{M-1}(\omega)$; by
Proposition~\ref{prop:fibers}(2) its span is $\mathrm{Hdg}_1$;
cyclic.
\item \emph{Weight $p \ge 2$:} the fiber $\mathrm{Hdg}_p = 0$
(Proposition~\ref{prop:fibers}(3)), and the zero orbit spans it;
cyclic.  (The honest shape: the certificate exists because the
designed fiber is dead, \emph{not} because of a vanishing theorem.)
\end{itemize}
(In Lean: \texttt{native\_GST\_single\_fold\_orbit\_law\_finiteLine\_unconditional\_gatefree},
wave 8-A3 \S4 --- the one-substitution kill shot.)
\end{proof}

\subsection{The consequences: the Hodge statement at the honest data}

\begin{definition}[The exact Hodge statement]
\label{def:exact-hodge}
$\mathrm{EveryHodgeClassIsRationalAlgebraic}(H)$ is the statement:
for every weight $p$ and every rational Hodge class $\alpha \in
\mathrm{Hdg}_p$, there is an actual codimension-$p$ cycle $Z$ of $V$
with
\[
\cl_p(Z) = \alpha .
\]
The finite rational combination form
$\mathrm{EveryHodgeClassIsFiniteRationalCombination}(H)$ is the
literal Clay wording: $\alpha$ is a finite $\Q$-combination of
classes of codimension-$p$ cycles.
\end{definition}

\begin{theorem}[One orbit law implies the Hodge statement]
\label{thm:consequence}
$\mathrm{NativeGSTSingleFoldOrbitLaw}(V,H)$ implies
$\mathrm{EveryHodgeClassIsRationalAlgebraic}(H)$, and equivalently
the finite rational combination form.
\end{theorem}

\begin{proof}
Let $\alpha \in \mathrm{Hdg}_p$ and let $(W,Z)$ be the cyclic
certificate at $p$.  Cyclicity gives $\alpha = \sum_{n} q_n W^n
\cl_p(Z)$, a finite rational combination.  Each orbit element is
algebraic: the seed class $\cl_p(Z)$ is the class of the actual cycle
$Z$, and the operator $W$ preserves algebraicity --- the
correspondence action on cycle classes is the push-pull along the
projections of the graph, and pushforwards and pullbacks of algebraic
cycles along algebraic correspondences are algebraic (the orbit
algebraicity receipt, the certificate's own \texttt{exists\_cycle}
method of \texttt{RealizedCorrespondenceKrylovCertificate}; realized on
points by Theorem~\ref{thm:realized} at the two live weights).  So
$W^n\cl_p(Z) = \cl_p(Z_n)$ for an actual cycle $Z_n$, and by
linearity of $\cl_p$:
\[
\alpha \;=\; \sum_n q_n \cl_p(Z_n) \;=\;
\cl_p\Bigl( \sum_n q_n Z_n \Bigr),
\]
the class of an actual codimension-$p$ cycle.  The finite combination
form is the same sum read term by term.  (In Lean:
\texttt{exactHodge\_of\_one\_GST\_fold\_orbit} /
\texttt{finiteRationalCombination\_of\_one\_GST\_fold\_orbit},
through \texttt{C.exists\_cycle} and the equivalence
\texttt{rationalAlgebraic\_iff\_finiteRationalCombination}.)
\end{proof}

\begin{corollary}[The Hodge conjecture at the honest data]
\label{cor:hodge-at-data}
For every $M \ge 1$ and every $\omega$: every rational Hodge class of
$H_M(\omega)$, of every weight, is the class of an actual
codimension-$p$ algebraic cycle of $V_M$, in the finite rational
combination form of the official wording.
\end{corollary}

\begin{proof}
Theorems~\ref{thm:law} and~\ref{thm:consequence}.
\end{proof}

% ======================================================================
\section{The universal primitive-X source}\label{sec:source}
% ======================================================================

\begin{definition}[The source families]\label{def:families}
The \emph{weight-zero source family} is the component family of
generic-point atoms:
\[
S^{(0)}_i \;=\; 1 \cdot \eta_i \qquad (i < M)
\qquad
(\texttt{finiteLineSourceZero}),
\]
and the \emph{weight-one source family} is the closed-point family
(the $R_1$-orbit of the honest seed):
\[
S^{(1)}_i \;=\; 1 \cdot \zeta_i \qquad (i < M)
\qquad
(\texttt{finiteLineSourceOne}).
\]
Their classes are, by Theorem~\ref{thm:seeds} (general $i$):
$\cl_0(S^{(0)}_i) = c_i$ and $\cl_1(S^{(1)}_i) = p_i(\omega)$.
\end{definition}

\begin{theorem}[Weight-zero independence; unconditional]
\label{thm:ind0}
The component family $(c_i)_{i<M}$ is $\Q$-linearly independent ---
unconditionally, at every family index.
\end{theorem}

\begin{proof}
Under the identification $H^0 \cong \mathrm{Fun}(M,\Q)$ the family is
the standard basis $(e_i)$: a dependence $\sum_i q_i c_i = 0$ reads
$\sum_i q_i e_i = 0$ in the function space, i.e. $q_i = 0$ for every
$i$ by evaluation.  (In Lean:
\texttt{finiteLineSourceZero\_independent}, through the equiv
readout and \texttt{EquivLike.injective}.)
\end{proof}

\begin{theorem}[Weight-one independence; at the family index]
\label{thm:ind1}
If $\omega \ne 0$, the point family $(p_i(\omega))_{i<M}$ is
$\Q$-linearly independent.  \emph{At $\omega = 0$ the family is
identically zero and the unguarded statement is false --- the fence is
exact.}
\end{theorem}

\begin{proof}
Suppose $\sum_i q_i\, p_i(\omega) = 0$ in
$H^2(\coprod_{j<M} S^2;\Q)$.  Restrict to the $j$-th slice: the
on-slice restriction kills every $p_i$ with $i \ne j$ (the collapse
$\circ$ inclusion is the constant-map composite off the slice) and
returns $\omega$ on $p_j$ by the section
property~\eqref{eq:section}.  So
\[
0 \;=\; \mathrm{incl}_j^* \Bigl( \sum_i q_i p_i(\omega) \Bigr) \;=\;
q_j\, \omega ,
\]
and $\omega \ne 0$ forces $q_j = 0$ --- for every $j$.  (In Lean:
\texttt{finiteLineSourceOne\_independent} at \texttt{hω};
the slice-restriction calculus is the green
\texttt{pullback\_pair} machinery.)
\end{proof}

\begin{theorem}[The spanning and the vacuity]
\label{thm:span}
Every weight-$0$ rational Hodge class is a $\Q$-combination of the
component family; every weight-$1$ rational Hodge class is a
$\Q$-combination of the closed-point family; every weight-$p \ge 2$
rational Hodge class is zero.
\end{theorem}

\begin{proof}
The three sectors of Proposition~\ref{prop:fibers}: at weight $0$,
$\mathrm{Hdg}_0 = H^0 = \Span_\Q(c_i)$ (the family is a basis by
Theorem~\ref{thm:ind0}); at weight $1$, $\mathrm{Hdg}_1 =
\Span_\Q(p_i(\omega))$ by the span-bigrading fiber receipt; at
weight $p \ge 2$, $\mathrm{Hdg}_p = 0$.  (In Lean:
\texttt{finiteLine\_universal\_primitive\_X\_source\_zero} /
\texttt{\_one} / \texttt{\_above}, the span memberships through
\texttt{Submodule.mem\_span\_range\_iff\_exists\_fun}.)
\end{proof}

\begin{construction}[The compilers]\label{con:compilers}
The \emph{weight-zero compiler} takes the explicit coordinates: a
class $\alpha \in H^0$ is compiled from its readout $E(\alpha) \in
\mathrm{Fun}(M,\Q)$ as
\[
\mathrm{cmp}_0(\alpha) \;=\; \sum_i E(\alpha)_i\, S^{(0)}_i
\qquad (\texttt{finiteLineSourceCompilerZero}),
\]
computable content, no choice.  The \emph{weight-one compiler} takes
the Classical choice of the span certificate of
Theorem~\ref{thm:span}:
\[
\mathrm{cmp}_1(\alpha) \;=\; \sum_i q_i(\alpha)\, S^{(1)}_i
\qquad (\texttt{finiteLineSourceCompilerOne}),
\]
noncomputable --- the corpus has no green degree-two dual-basis
machine; the expected \texttt{Classical.choice} axiom receipt is the
designated honesty marker.
\end{construction}

\begin{theorem}[Compiler correctness]
\label{thm:compilers}
$\cl_0 \circ \mathrm{cmp}_0 = \mathrm{id}$ on all of $H^0$, and
$\cl_1 \circ \mathrm{cmp}_1 = \mathrm{id}$ on
$\mathrm{Hdg}_1$.
\end{theorem}

\begin{proof}
Weight zero: $\cl_0(\sum_i E(\alpha)_i S^{(0)}_i) = \sum_i
E(\alpha)_i c_i = \alpha$ because $(c_i)$ is the standard basis
under the readout $E$.  Weight one: the chosen coordinates realize
$\alpha$ in the span by Theorem~\ref{thm:span}.  (In Lean:
\texttt{finiteLineSourceCompilerZero\_correct} /
\texttt{finiteLineSourceCompilerOne\_correct}.)
\end{proof}

\begin{theorem}[The universal primitive-X source package; the
$\forall\omega$ master]\label{thm:master}
For every $M \ge 1$ and every family index $\omega \in H^2(S^2;\Q)$,
at the honest data $H_M(\omega)$: the weight-zero family is
independent; the weight-one family is independent whenever $\omega
\ne 0$; both live sectors are spanned by their families; the weight
$p \ge 2$ sector is zero; both compilers are correct; and the
single-fold orbit law holds.
\end{theorem}

\begin{proof}
The conjunction of Theorems~\ref{thm:ind0}, \ref{thm:ind1},
\ref{thm:span}, \ref{thm:compilers}, and \ref{thm:law}.  (In Lean:
\texttt{the\_universal\_primitive\_X\_source\_package}, wave 9
green, run 45736114591 at \texttt{5c98f94} --- the eight conjuncts
assembled.)
\end{proof}

% ======================================================================
\section{The fence: the tetrahedron machine and the exists-flip}
\label{sec:fence}
% ======================================================================

The single hypothesis the master package carries beyond the corpus is
the \emph{nonzero family index}.  Wave 10 built the machinery to flip
the quantifier: \emph{given one} nonzero degree-two sphere class, the
package instantiates at that witness.

\subsection{The tetrahedron machine}

\begin{construction}[The projected tetrahedron]
\label{con:tetra}
In $\R^3$ with the Euclidean norm:
\begin{itemize}
\item the regular tetrahedron $T = \mathrm{conv}\{w_0, w_1, w_2,
w_3\}$ with
\[
w_0 = (1,1,1),\quad w_1 = (1,-1,-1),\quad w_2 = (-1,1,-1),\quad
w_3 = (-1,-1,1);
\]
\item the camera $c = (10,10,10)$;
\item the central projection
\[
F : \R^3 \setminus \{0\} \to S^2, \qquad F(y) = \|y\|^{-1} y ,
\]
continuous at every nonzero point (the scalar factor $y \mapsto
\|y\|^{-1}$ is continuous away from $0$ and is multiplied against the
identity);
\item the barycentric parametrization $\beta : \Delta^3 \to T$,
$\beta(x) = \sum_i x_i w_i$ on the standard $3$-simplex;
\item the singular $3$-simplex of the projected tetrahedron:
\[
\mathrm{tet} : \Delta^3 \to S^2, \qquad \mathrm{tet}(x) = F\bigl(
\beta(x) - c \bigr),
\]
continuous since $\beta(x) - c \ne 0$ always
(Lemma~\ref{lem:camera} below);
\item the generator $\mathrm{tetGen} \in C_3(S^2;\Q)$, the singular
$3$-chain carried by $\mathrm{tet}$; and the \emph{projected
boundary cycle}
\[
z_{\mathrm{tet}} \;=\; \partial_3(\mathrm{tetGen}) \;\in\;
C_2(S^2;\Q),
\]
the alternating sum of the four projected $2$-faces.
\end{itemize}
\end{construction}

\begin{lemma}[The camera separation]\label{lem:camera}
For every $x \in \Delta^3$: $\beta(x) \ne c$.  Indeed the first
coordinate of any point of $T$ is at most $1$, while the camera's
first coordinate is $10$.
\end{lemma}

\begin{proof}
Every vertex first-coordinate is $\pm 1 \le 1$; a convex combination
$\sum_i x_i (w_i)_1$ with $x_i \ge 0$, $\sum_i x_i = 1$ satisfies
\[
\sum_i x_i (w_i)_1 \;\le\; \sum_i x_i \cdot 1 \;=\; 1 \;<\; 10 .
\]
(In Lean: \texttt{tetraBary\_ne\_camera}, the upper-bound route ---
\texttt{mul\_le\_mul\_of\_nonneg\_left} + \texttt{Finset.sum\_le\_sum}
+ the simplex sum property; no absolute value of any sum appears.)
\end{proof}

\begin{theorem}[The cycle law and the face distinctness]
\label{thm:tetra}
\begin{enumerate}
\item $\partial_2 z_{\mathrm{tet}} = 0$ --- the projected boundary is
a cycle, by the complex law $\partial \circ \partial = 0$;
\item the four projected faces are pairwise distinct as singular
simplices; in particular the $0$-th and $1$-st faces differ, by the
two-coordinate contradiction: if
$F(w_0 - c) = F(w_1 - c)$, then with $r_i = \|w_i - c\|$:
\[
\underbrace{r_0^{-1} \cdot (-9) = r_1^{-1} \cdot (-9)}_{\text{coordinate
}0} \;\wedge\;
\underbrace{r_0^{-1} \cdot (-9) = r_1^{-1} \cdot (-11)}_{\text{coordinate
}1}
\;\Longrightarrow\; r_0 = r_1 \;\Longrightarrow\; -9 = -11,
\]
since $w_0 - c = (-9,-9,-9)$ and $w_1 - c = (-9,-11,-11)$.
\end{enumerate}
\end{theorem}

\begin{proof}
(1) is the pinned chain-complex law applied to
$\mathrm{tetGen}$.  (2): coordinate $0$ of the equality
$r_0^{-1}(w_0 - c) = r_1^{-1}(w_1 - c)$ reads $r_0^{-1}(-9) =
r_1^{-1}(-9)$; cancelling $-9 \ne 0$ gives $r_0^{-1} = r_1^{-1}$;
coordinate $1$ then reads $r_0^{-1}(-9) = r_0^{-1}(-11)$, and
cancelling $r_0^{-1} \ne 0$ gives $-9 = -11$, absurd.  The vertex
evaluations ride \texttt{stdSimplex.map\_vertex} +
\texttt{tetraBary\_vertex}.  (In Lean: \texttt{tetraCycle\_d} and
\texttt{sphereNormalize\_w0\_ne\_w1}, wave 10-B green, run
45780735883 at \texttt{00463dc}.)
\end{proof}

\subsection{The separator shape and the exists-flip master}

\begin{definition}[The separator shape]\label{def:separator}
$\mathrm{TetraSeparator}$ is the proposition: there exists a rational
linear functional $\varphi : C_2(S^2;\Q) \to \Q$ that vanishes on
every degree-three boundary and reads $1$ on the projected boundary
cycle:
\[
\exists \varphi : C_2(S^2;\Q) \to \Q, \quad
\bigl( \forall w \in C_3,\; \varphi(\partial_3 w) = 0 \bigr)
\;\wedge\; \varphi(z_{\mathrm{tet}}) = 1 .
\]
\end{definition}

The realization design (the classical argument, recorded in the
module header of \texttt{GSTSphereNonzeroClass.lean}): a separator is
a cocycle of the rational dual complex (the dual differential is
precomposition with $\partial_3$, so boundary-vanishing \emph{is} the
cocycle law); its class $\omega_0 = [\varphi] \in H^2(S^2;\Q)$ is
nonzero, because a coboundary $\delta\psi = \psi \circ \partial_2$
vanishes on every cycle and so cannot read $1$ on
$z_{\mathrm{tet}}$.

\begin{theorem}[The exists-flip master]\label{thm:exists-flip}
Given one nonzero degree-two sphere class --- the fence hypothesis
\[
(\Omega):\quad \exists\, \omega \in H^2(S^2;\Q),\; \omega \ne 0
\]
--- the full package instantiates at that witness: there exists
$\omega$ such that the weight-one family is independent at every $M
\ge 1$, the weight-zero family is independent, both live sectors are
spanned, the $p \ge 2$ sector is zero, both compilers are correct,
and the gate-free law holds.  In symbols:
\[
(\Omega) \;\Longrightarrow\; \exists \omega,\;
\text{the eight conjuncts of Theorem~\ref{thm:master} at }
\omega .
\]
\end{theorem}

\begin{proof}
Take the witness $\omega_0$ of $(\Omega)$; each conjunct is the
corresponding $\forall\omega$ theorem of the certified corpus
instantiated at $\omega_0$ (the independence at $\omega_0 \ne 0$;
the rest are $\forall\omega$).  The proof is a witness
instantuation: no new mathematics, only the quantifier flip the
master statement of record demands.  (In Lean:
\texttt{finiteLineSourceOne\_independent\_of\_exists} and
\texttt{the\_unconditional\_source\_package\_of\_exists}, wave 10-A
green, run 45780735883 at \texttt{00463dc}.)
\end{proof}

% ======================================================================
\section{The route correction: the projected boundary bounds}
\label{sec:correction}
% ======================================================================

The wave-10 design designated, as the next append of the fence lane,
the construction of a separator \emph{from} the projected boundary
cycle (the FENCE-Z route: prove $[z_{\mathrm{tet}}] \ne 0$ in
$H_2(S^2;\Q)$, then build the separating functional on the
homology).  The route is \emph{closed} by the following kernel
certificate, pushed this round.

\begin{theorem}[The projected boundary bounds its own generator;
wave 10-D]\label{thm:bounds}
\[
\partial_3(\mathrm{tetGen}) \;=\; z_{\mathrm{tet}}
\qquad\text{by construction, hence}\qquad
[z_{\mathrm{tet}}] \;=\; 0 \;\in\; H_2(S^2;\Q).
\]
\end{theorem}

\begin{proof}
The cycle was \emph{defined} as the boundary of its own generator:
$z_{\mathrm{tet}} := \partial_3(\mathrm{tetGen})$
(Construction~\ref{con:tetra}).  The generator is an honest
degree-three singular chain of the sphere --- the singular
$3$-simplex $\mathrm{tet} : \Delta^3 \to S^2$ is a continuous map
from the solid tetrahedron \emph{into} the sphere, defined and
continuous on all of $\Delta^3$ precisely because the camera sits
outside $T$ (Lemma~\ref{lem:camera}).  The boundary of any singular
chain is null-homologous --- it bounds that chain.  (In Lean:
\texttt{tetraCycle\_bounds} --- \texttt{sphereBoundaryThree tetraGen
= tetraCycle} by unfolding both definitions --- and
\texttt{tetraCycle\_is\_boundary}, the explicit filling witness;
wave 10-D, this round.)
\end{proof}

\begin{theorem}[The separator shape is false at this cycle]
\label{thm:not-separator}
$\neg\,\mathrm{TetraSeparator}$: no rational functional vanishing on
all degree-three boundaries can read $1$ on
$z_{\mathrm{tet}}$.
\end{theorem}

\begin{proof}
Let $\varphi$ vanish on all boundaries.  Then
\[
\varphi(z_{\mathrm{tet}}) \;=\;
\varphi\bigl( \partial_3(\mathrm{tetGen}) \bigr) \;=\; 0
\]
by Theorem~\ref{thm:bounds} and the vanishing, while the separator
demand reads $\varphi(z_{\mathrm{tet}}) = 1$; and $0 \ne 1$.  (In
Lean: \texttt{not\_TetraSeparator} --- three lines:
\texttt{rintro}, the rewrite by \texttt{tetraCycle\_bounds},
\texttt{zero\_ne\_one}; and the positive form
\texttt{separator\_vanishes\_on\_tetraCycle}.)
\end{proof}

\begin{remark}[What died and what survived]
\begin{itemize}
\item \textbf{Died:} the FENCE-Z route at this cycle --- the
designated next appends TET-6 (the nonvanishing theorem at
$z_{\mathrm{tet}}$) and NZC-5 (the separator construction from
$[z_{\mathrm{tet}}] \ne 0$) are cancelled: their premise is false.
The mathematical reason is structural, not repairable: the projected
\emph{solid} tetrahedron is a singular $3$-chain \emph{of the
sphere} (every continuous map into the sphere is), so its boundary
is null-homologous by construction.  No single-simplex projection
can ever produce a non-bounding cycle this way.
\item \textbf{Survived untouched:} the entire certified corpus.  The
wave-10 trio's green statements (the cycle law, the camera
separation, the face distinctness, the separator \emph{shape as a
definition}) are all true; the exists-flip master takes $(\Omega)$
as an explicit hypothesis and never used the route; the seven tasks
of \S\ref{sec:world}--\ref{sec:source} are independent of it.
\item \textbf{The corrected route:} a genuinely non-bounding cycle
is required --- the classical generator of $H_2(S^2;\Q) \cong \Q$,
constructed and proved in \S\ref{sec:omega}.  Its Lean
formalization needs sphere-topology infrastructure (a
Mayer-Vietoris sequence, degree theory, or a no-retraction theorem)
that the pinned Mathlib does not carry --- verified directly against
the pin: \texttt{Mathlib/AlgebraicTopology/} at \texttt{51e6992e}
contains Dold-Kan, simplicial sets, the singular chain complex,
$H_0$, homotopy invariance, and the Eilenberg-Steenrod
\emph{axioms}, but no sphere computation.  This is the campaign's
designated research front, now with the poisoned route closed by
certificate instead of burning CI rounds against an impossible
target.
\end{itemize}
\end{remark}

% ======================================================================
\section{The classical completion: \texorpdfstring{$H^2(S^2;\Q) \cong
\Q$}{H2(S2;Q) = Q}}\label{sec:omega}
% ======================================================================

This section proves the fence hypothesis $(\Omega)$ classically, in
full, by a self-contained computation.  Together with
Theorem~\ref{thm:exists-flip} this closes the classical form of the
main theorem: \emph{the entire package of
Theorem~\ref{thm:master} holds unconditionally at every $M \ge 1$
and every nonzero $\omega$, and such $\omega$ exist.}

\begin{theorem}[The fence]\label{thm:sphere-h2}
$H^2(S^2;\Q) \ne 0$.  In fact $H_2(S^2;\Q) \cong \Q$ and, $\Q$
being a field, $H^2(S^2;\Q) \cong \Hom_\Q(H_2(S^2;\Q), \Q) \cong \Q$
by the universal coefficient theorem.
\end{theorem}

\begin{proof}
\emph{Step 1: the simplicial model.}  Let $K = \partial\Delta^3$ be
the boundary of the standard $3$-simplex --- the abstract simplicial
complex on vertices $\{0,1,2,3\}$ whose simplices are the proper
nonempty subsets.  Its geometric realization $|K|$ is the boundary of
the tetrahedron, and the radial projection from the barycenter
\[
\rho : |K| \to S^2, \qquad y \mapsto y/\|y\|
\]
is a homeomorphism (the barycenter is an interior point of the
tetrahedron, so $\rho$ is a well-defined continuous bijection from a
compact space to a Hausdorff space).  Singular and simplicial
homology of $|K|$ agree (the standard theorem: the simplicial chain
complex computes singular homology of the realization; over $\Q$
either via the subdivision operator or the $\Delta$-complex
comparison).

\emph{Step 2: the chain groups.}  Orient each face by increasing
vertex order and write $F_i$ for the face opposite vertex $i$:
\[
F_0 = [1,2,3],\; F_1 = [0,2,3],\; F_2 = [0,1,3],\; F_3 = [0,1,2].
\]
The simplicial chain groups of $K$ are
\[
C_2 = \Q^4\langle F_0,F_1,F_2,F_3\rangle, \qquad
C_1 = \Q^6\langle E_{jk} : j < k \rangle, \qquad
C_0 = \Q^4, \qquad C_3 = 0
\]
($K$ has no $3$-simplices).  The face boundaries, from
$\partial[v_0v_1v_2] = [v_1v_2] - [v_0v_2] + [v_0v_1]$:
\[
\begin{aligned}
\partial F_0 &= E_{23} - E_{13} + E_{12}, &
\partial F_1 &= E_{23} - E_{03} + E_{02},\\
\partial F_2 &= E_{13} - E_{03} + E_{01}, &
\partial F_3 &= E_{12} - E_{02} + E_{01} .
\end{aligned}
\]

\emph{Step 3: the kernel of $\partial_2$, explicitly.}  Let $a =
(a_0,a_1,a_2,a_3)$ with $\sum_i a_i \partial F_i = 0$.  Collecting
edge coefficients:
\[
\begin{array}{llll}
E_{01}: & a_2 + a_3 = 0, \qquad & E_{02}: & a_1 - a_3 = 0,\\
E_{03}: & a_1 + a_2 = 0, \qquad & E_{12}: & a_0 + a_3 = 0,\\
E_{13}: & a_0 - a_2 = 0, \qquad & E_{23}: & a_0 + a_1 = 0 .
\end{array}
\]
The first, fourth and fifth equations give $a_2 = -a_3$, $a_0 =
-a_3$, $a_0 = a_2$ (consistent); the second and third give $a_1 =
a_3$, $a_1 = -a_2 = a_3$ (consistent); the sixth, $a_0 + a_1 = -a_3
+ a_3 = 0$, is then automatic.  So
\[
\ker \partial_2 \;=\; \Q \cdot (1, -1, 1, -1),
\]
one-dimensional, generated by the fundamental cycle
\[
z_{\mathrm{fund}} \;=\; F_0 - F_1 + F_2 - F_3 \;=\;
\partial [0,1,2,3],
\]
the boundary of the oriented solid simplex --- a cycle by $\partial
\circ \partial = 0$, and by the computation every cycle is a
multiple of it.

\emph{Step 4: the conclusion.}  Since $C_3(K) = 0$ --- there are no
$3$-simplices in the boundary complex --- there are no boundaries to
kill:
\[
H_2(K;\Q) \;=\; \ker \partial_2 / \operatorname{im} \partial_3
\;=\; \ker \partial_2 \;\cong\; \Q .
\]
Transporting along the homeomorphism $\rho$:
$H_2(S^2;\Q) \cong \Q \ne 0$, and by the universal coefficient
theorem over the field $\Q$ (no $\operatorname{Ext}$ terms):
$H^2(S^2;\Q) \cong \Hom_\Q(\Q, \Q) \cong \Q \ne 0$.
\end{proof}

\begin{remark}[Why the solid tetrahedron does not contradict this]
The fundamental cycle $z_{\mathrm{fund}}$ of Step 3 is the boundary
of the solid simplex \emph{which is not a chain of $K$} --- the
boundary complex has no $3$-simplices, so within $K$ the cycle
cannot bound.  This is exactly the structural difference the route
correction of \S\ref{sec:correction} exposed: the projected solid
tetrahedron \emph{is} a chain of the sphere (the map lands inside
the space), so there the boundary \emph{does} bound.  A non-bounding
cycle on $S^2$ must come from a triangulation \emph{of the sphere
itself}, not from a projection of a solid body into it.
\end{remark}

\begin{construction}[The classical witness and separator]
\label{con:classical-separator}
Theorem~\ref{thm:sphere-h2} supplies the witness: any generator
$\omega$ of $H^2(S^2;\Q) \cong \Q$, concretely the class dual to
$z_{\mathrm{fund}}$ under the transport.  Equivalently, the
separator route closed in \S\ref{sec:correction} for the projected
cycle opens classically for the fundamental cycle: the coordinate
functional $\mathcal{L} : H_2(S^2;\Q) \to \Q$ with
$\mathcal{L}([z_{\mathrm{fund}}]) = 1$ composes with the homology
projection and extends from the cycles submodule to all of $C_2$
(vector-space extension), yielding a functional vanishing on
boundaries and reading $1$ on a \emph{non-bounding} cycle --- a
genuine separator, whose realization (the cocycle-to-class argument
of \S\ref{sec:fence}) is the nonzero class $\omega$.
\end{construction}

\begin{theorem}[The closed package; classical form]
\label{thm:closed}
The fence hypothesis $(\Omega)$ holds
(Theorem~\ref{thm:sphere-h2}); therefore, by
Theorem~\ref{thm:exists-flip}, there exists $\omega \in
H^2(S^2;\Q)$, $\omega \ne 0$, such that at every $M \ge 1$ the
honest data $H_M(\omega)$ carries the full eight-conjunct package of
Theorem~\ref{thm:master}: both source families bases of their
sectors, both compilers correct, the higher sector zero, and the
single-fold orbit law --- the Hodge conjecture statement at the
honest data, with explicit witnesses.  This is
Theorem~\ref{thm:main}.
\end{theorem}

\begin{remark}[The formalization boundary, stated honestly]
The Lean corpus certifies every ingredient of
Theorem~\ref{thm:closed} except the fence
(Theorem~\ref{thm:sphere-h2}), which enters the machine-checked
master as an explicit hypothesis.  The gap is infrastructure, not
mathematics: no Mayer-Vietoris sequence, no degree theory, no
no-retraction theorem, and no sphere homology exist in the pinned
Mathlib, and the wave-10 Lean route to a witness was closed by the
route-correction certificate
(Theorem~\ref{thm:not-separator}).  The formalization of
Theorem~\ref{thm:sphere-h2} --- e.g.\ via the $\partial\Delta^3$
simplicial computation plus the realization-to-singular comparison,
or via a no-retraction theorem --- is the campaign's designated
research front.  No green is claimed beyond what CI certified at the
pushed SHAs.
\end{remark}

\subsection{The formal correction begun: the center-radial cycle in
Lean}

The boundary above stated the front as of wave 10.  Wave 11 (the
module \texttt{GSTSphereFundamentalCycle.lean}, the fence-discharge
lane's fifth ladder piece) constructs the classical generator itself
in the honest singular chain complex of the sphere --- the corrected
route's first certified half:

\begin{construction}[The fundamental cycle, formalized]
\label{con:fund-cycle-lean}
The four vertices of the alternating tetrahedron (the corpus's
\texttt{tetraW}) have the origin as barycenter --- the sphere's own
center.  For each face $F_i$ (the three vertices other than $i$, in
increasing order), the center-radial projection $y \mapsto
\|y\|^{-1} y$ is well defined on the whole closed face: any three of
the four vertices are linearly independent (each face's coordinate
matrix is a $\pm 1$ matrix with determinant $\pm 4$), so no closed
face meets the origin --- while the solid does, which is exactly why
the projected boundary of the \emph{camera} tetrahedron bounds
(\S\ref{sec:correction}) and this cycle does not.  Let $\sigma_i :
\Delta^2 \to S^2$ be the composition (a singular $2$-simplex of the
sphere) and
\[
\texttt{fundCycle} \;=\; \iota(\sigma_0) - \iota(\sigma_1) +
  \iota(\sigma_2) - \iota(\sigma_3) \;\in\;
  C_2^{\mathrm{sing}}(S^2;\Q)
\]
the alternating sum of the injected face chains.
\end{construction}

\begin{theorem}[The cycle law, machine-checked]
\label{thm:fund-cycle-d}
$d_2(\texttt{fundCycle}) = 0$ (\texttt{fundCycle\_d}), by the
twelve-term edge cancellation: each of the six tetrahedron edges
appears in exactly two faces with opposite total sign and the
\emph{same} parametrized singular $1$-simplex --- the affine coface
calculus (the pushed-forward point's coordinates are the fiber sums
of the pinned \texttt{stdSimplex.map}, and both visits of an edge
compose to the same affine parametrization of the same vertex pair)
--- so the six pairs cancel term by term.
\end{theorem}

\begin{construction}[The corrected separator shape]
\label{con:fund-separator}
The route correction (\S\ref{sec:correction}) proved the separator
shape false at the camera cycle.  Restated at the fundamental cycle:
\[
\texttt{FundSeparator} \;:\;\;
  \exists\, \varphi : C_2^{\mathrm{sing}}(S^2;\Q) \to_\Q \Q,\;
  \bigl(\forall w,\ \varphi(d_3 w) = 0\bigr) \;\wedge\;
  \varphi(\texttt{fundCycle}) = 1 ,
\]
with the consistency receipt \texttt{no\_separator\_of\_boundary}: if
the fundamental cycle bounded, the shape would be false --- the shape
lives exactly as long as the cycle is non-bounding.  Classically
(Theorem~\ref{thm:sphere-h2}) the cycle is non-bounding, so the shape
holds (extend the degree-one functional on the line it spans, modulo
the boundary submodule, to all of $C_2$).
\end{construction}

The remaining formalization step --- frozen as the campaign's next
wave, not claimed here --- is the realization: every separator is a
cocycle of the dual cochain complex (the boundary-vanishing is the
cocycle law), and its class is nonzero, since a coboundary $\psi
\circ d_2$ reads $\psi(d_2\,\texttt{fundCycle}) = \psi(0) = 0$ on the
cycle by Theorem~\ref{thm:fund-cycle-d} and so cannot be the
separator that reads $1$.  That single step discharges $(\Omega)$;
with the exists-flip master it then discharges the two stashed
sectors of the honest data (the weight-one spanning at a nonzero
family index, the higher-weight vanishing), completing the
topological upgrade of the point arithmetic.

% ======================================================================
\section{The certification receipts}\label{sec:certification}
% ======================================================================

Every theorem above is machine-checked.  The receipts, at the branch
\texttt{research/gpt6-gst-hypothesis-part1-chow-square-20261110} of
\texttt{kyo-oo/HC-PROOF} (Lean 4.33.0-rc2, Mathlib pin
\texttt{51e6992efd06126df61a496bebf8f49482a4e129}):

\begin{longtable}{@{}p{0.30\textwidth}p{0.34\textwidth}p{0.28\textwidth}@{}}
\toprule
\textbf{Component (this paper)} & \textbf{Lean module} &
\textbf{CI lane / verdict}\\
\midrule
\endhead
World, ring, embedding (\S\ref{sec:world}) &
\texttt{GSTFiniteLineRing}, \texttt{GSTFiniteLineOverBase},
\texttt{GSTFiniteLineWorldCarrier} & green, unconditional\\
Prime lattice, retractions, no-intermediate
(Thms~\ref{thm:kernel}--\ref{thm:no-intermediate}) &
\texttt{GSTFiniteLinePointAlgebra} &
\texttt{gst-finite-line-point-algebra}: green\\
Stereographic equivalence (Thm~\ref{thm:stereo}) &
\texttt{GSTFiniteLineRiemann} & \texttt{gst-finite-line-riemann}:
green\\
Analytification, charts (\S\ref{sec:analytification}) &
\texttt{GSTFiniteLineAnalytification} &
\texttt{gst-finite-line-analytification}: green\\
Points gate, GATE-A (Thm~\ref{thm:gate}) &
\texttt{GSTFiniteLinePointsGateCore} + \texttt{Discharge} (+ InverseA/B) &
\texttt{gst-finite-line-points-gate-discharge}: green\\
Honest Hodge data, seeds (Thms~\ref{thm:seeds},
\ref{prop:fibers}) & \texttt{GSTFiniteLineHonestHodgeData} &
\texttt{gst-finite-line-honest-hodge-data}: green\\
The two letters (Thms~\ref{thm:realized}) &
\texttt{GSTFiniteLineCycleCorrespondence} (+ Basic, Action,
Operators) & \texttt{gst-finite-line-cycle-correspondence}:
green\\
The law, gate-free (Thm~\ref{thm:law}) &
\texttt{GSTFiniteLineUnconditionalSource},
\texttt{GSTFiniteLinePointsGateDischarge} \S4 &
\texttt{gst-finite-line-unconditional-source}: green (run
45736114591 at \texttt{5c98f94} apex)\\
Consequences (Thm~\ref{thm:consequence}) &
\texttt{GSTHodgeSingleHypothesisConditional} & green, conditional
chain\\
Source families, independence, compilers
(Thms~\ref{thm:ind0}--\ref{thm:compilers}) &
\texttt{GSTFiniteLinePrimitiveXSourceFamilies},
\texttt{GSTFiniteLinePrimitiveXSource} &
\texttt{gst-finite-line-primitive-x-source}: green\\
Master package (Thm~\ref{thm:master}) &
\texttt{GSTFiniteLinePrimitiveXSource} & wave 9 green, run
45736114591 at \texttt{5c98f94}\\
Tetrahedron machine (Thms~\ref{thm:tetra}) &
\texttt{GSTSphereTetrahedronCycle} &
\texttt{gst-sphere-fence-discharge} Ladder 1/3: green at
\texttt{00463dc} (run 45780735883)\\
Separator shape (Def~\ref{def:separator}) &
\texttt{GSTSphereNonzeroClass} & Ladder 2/3: green at
\texttt{00463dc}\\
Exists-flip master (Thm~\ref{thm:exists-flip}) &
\texttt{GSTFiniteLineFenceDischarge} & Ladder 3/3: green at
\texttt{00463dc}\\
Route correction (Thms~\ref{thm:bounds},
\ref{thm:not-separator}) & \texttt{GSTSphereRouteCorrection} &
Ladder 4/4: this round's push (\S\ref{sec:correction})\\
Fence, classical (Thm~\ref{thm:sphere-h2}) & --- (classical
proof, \S\ref{sec:omega}) & carried as the explicit hypothesis
$(\Omega)$\\
\bottomrule
\end{longtable}

\noindent All lanes run the kernel check with a
\texttt{sorryAx}/\texttt{declaration uses 'sorry'} audit; the apex
runs additionally print \texttt{\#print axioms} receipts.  The honest
disclosure ledger of the corpus: (i) the weight-one independence
carries \texttt{h\ensuremath{\omega}~:~\ensuremath{\omega}~\ensuremath{\ne}~0} exactly as
Theorem~\ref{thm:ind1} states --- at $\omega = 0$ the family is
zero; (ii) the $p \ge 2$ vacuity is the bigrading design, not a
vanishing theorem; (iii) the weight-one compiler is Classical-choice
noncomputable; (iv) the package's law conjunct is stated at the
constructed unconditional analytification, never at an arbitrary
datum.

\subsection{Postscript: the discipline}

The campaign's standing law is that no green is claimed beyond what
CI certified at the pushed SHAs, and no statement is strengthened
beyond its frozen form.  This document adds the complementary
discipline: every green label is backed here by a complete
mathematical proof on paper.  Where the paper's mathematics exceeds
the certified corpus --- exactly one statement, the fence of
\S\ref{sec:omega} --- the excess is labeled as classical, its proof
is given in full, and its formalization front is designated.

\end{document}
